\section{The CLAIMFORGE Benchmark}
\label{sec:benchmark}

\subsection{Task Definition}

Each CLAIMFORGE example starts from a real source image $x$ and a human-marked insertion region $m$. An editor generates a modified context crop in which a claim-relevant object is inserted inside $m$. The benchmark then composes a forged full image $\tilde{x}$ by replacing only the selected object region or marked edit region in the source image. The released target contains two labels: an image-level authenticity label $y \in \{0,1\}$ and a binary pixel mask $M$ for the forged region.

The benchmark has two mandatory tasks:
\begin{description}
    \item[T1: Detection.] Given an image, predict whether it is real or forged. We report AUC, AP, and fixed-threshold accuracy.
    \item[T2: Localization.] Given a forged image, predict a pixel-level mask for the inserted region. We report F1, IoU, MCC, and pixel AP.
\end{description}
Methods that natively solve only one task can still participate through documented adapters. For example, a whole-image detector can provide coarse localization through patch scores or class-activation maps, but the result must be marked as a non-native localizer.

\subsection{Current Data Construction}

The current repository contains an expanded source pool and an active generation handoff. The source pool has 300 restaurant candidates and 300 lodging candidates screened from Open Images and replacement sources. The active handoff contains 297 complete restaurant edit slots. Each slot specifies a source image, a context crop, an insertion box, a full-resolution mask, and a prompt target. The current target object is ``mouse'' for the active restaurant batch.

\begin{table}[t]
\centering
\small
\begin{tabular}{lr}
\toprule
Repository quantity & Value \\
\midrule
Candidate source images & 600 \\
Restaurant candidates & 300 \\
Lodging candidates & 300 \\
Complete restaurant edit slots & 297 \\
Pilot HunyuanImage-3 generations & 22 \\
Median insertion width & 72 px \\
Median insertion height & 51 px \\
Median insertion area & 3,596 px \\
Median insertion/source area & 0.232\% \\
Mean insertion/source area & 0.363\% \\
\bottomrule
\end{tabular}
\caption{Current CLAIMFORGE repository state used by this draft. These are construction statistics, not detector results.}
\label{tab:repo-state}
\end{table}

\subsection{Pixel-Preserving Composition}

The construction pipeline is designed to isolate the inserted object from unrelated image changes. A generation agent receives the context crop and the crop-local edit coordinates. The prompt asks for a single small realistic object inside the target region while preserving lighting, shadows, texture, perspective, and JPEG-like realism. The composer then pastes the generated object back into the original source image. In the current pilot path, the composer can estimate an object-only mask by thresholding crop differences inside the insertion box, keeping the largest connected component, dilating it for contact shadows, and feathering the mask before compositing.

This design supports three controls that are essential for a clean benchmark claim. First, matched real and forged images should share the same final post-processing so that format and compression cues cannot separate the classes. Second, a real-object paste-back control should use the same paste pipeline with a real cutout, disentangling AI texture from paste artifacts. Third, real images containing genuine claim objects should be included so that a model cannot solve the task by detecting the presence of a pest, stain, or damage object.

\subsection{Scale-Up Matrix}

The current handoff is a construction snapshot, not the final benchmark scale. The intended full benchmark expands across domains, editing models, object categories, and laundering transformations. Domains include restaurant, room, bathroom, kitchen, short-term rental, and related claim contexts. Editors should include HunyuanImage-3 plus several open and commercial inpainting/editing systems. Object categories should include pests, contamination, and damage. Laundering includes JPEG quality sweeps, resizing, blur/noise, social-media round trips, screenshots, and double compression. The final split should hold out at least one editor and one domain for zero-shot generalization.
